%% ==========================================================================
%%  Appendix C  The algorithm for objective signed dichotomous valuations
%%  (from the supplementary material of the AAMAS 2027 submission, Section 3)
%% ==========================================================================
\section{Algorithm for Objective Signed Dichotomous Valuations}
\label{app:objective-mixed-alg}
We first restate the theorem for convenience.
\renewcommand{\restatedname}{Theorem~\ref{thm:objective-mixed}}
\begin{restated}
For every instance with objective signed dichotomous valuations, there exist a complete allocation $\alloc$ and a subsidy vector $p\in\{0,1\}^n$ such that $(\alloc,p)$ is envy-free and
\[
\sum_{i\in N} p_i\leq n-1.
\]
Moreover, such an allocation and subsidy vector can be computed in
polynomial time in the value-oracle model.
\end{restated}
We next give the complete algorithm underlying Theorem~\ref{thm:objective-mixed}, using Lemma~\ref{lem:goods-extension} as the one-good extension subroutine.

For brevity, let \textsc{ChoreAllocate} denote the polynomial-time procedure developed in Section~\ref{sec:main} (Algorithm~\ref{alg:main}) for negative dichotomous valuations. Given a negative dichotomous chore instance, it returns a complete allocation $\alloc$ and $p\in\{0,1\}^n$ such that $(\alloc,p)$ is envy-free and $\sum_i p_i\le n-1$.

\begin{algorithm}[H]
\caption{Unit subsidies for objective signed dichotomous valuations}
\label{alg:objective-mixed}
\begin{algorithmic}[1]
\Require Agents $N=[n]$, items $M=G\mathbin{\dot\cup}C$, and value-oracle access to objective signed dichotomous valuations $\{v_i\}_{i\in N}$.
\Ensure A complete allocation $\alloc$ and $p\in\{0,1\}^n$ such that $(\alloc,p)$ is envy-free and $\sum_{i\in N}p_i\le n-1$.

\State Restrict each $v_i$ to subsets of $C$.
\State $(\alloc,p)\gets\textsc{ChoreAllocate}(N,C,\{v_i|_C\}_{i\in N})$.
\State Replace $p$ by the pointwise-minimal subsidy vector supporting $\alloc$.

\For{each $g\in G$}
    \State Apply the one-good extension procedure of Lemma~\ref{lem:goods-extension} to $(\alloc,p)$ and $g$, obtaining an allocation $\alloc'$ of the currently allocated items together with $g$ and a subsidy vector $p'\in\{0,1\}^n$ such that $(\alloc',p')$ is envy-free.
    \State $\alloc\gets\alloc'$.
    \State Replace $p'$ by the pointwise-minimal subsidy vector supporting $\alloc$, and set $p\gets p'$.
\EndFor

\State \Return $(\alloc,p)$.
\end{algorithmic}
\end{algorithm}

The one-good extension in Line~5 is implemented by the \textsc{Extend}/\textsc{FindSink} procedure of Barman et al.~\cite{BKNS22}, whose validity in the present setting is established in the proof of Lemma~\ref{lem:goods-extension} (Appendix~\ref{app:goods-extension}). After every iteration, we take the pointwise-minimal supporting subsidy vector. Since the valuations are integer-valued, this vector is integral, and Lemma~\ref{lem:goods-extension} implies that all its components are at most one. Moreover, pointwise minimality implies that at least one component is zero. Hence, throughout the algorithm
\[
p\in\{0,1\}^n
\qquad\text{and}\qquad
\sum_{i\in N}p_i\le n-1.
\]
